Let $X$ be the connected sum of three real projective planes.
%%%%%%%%%%%%%%%%%%%%%%%%%%%%%%%%%%%%%%%%%%%%%%%%%%%%%%%%%%%%%%%%%%%
\begin{alist}
%------------------------------------------------------------------*
\item Describe a CW decomposition of $X$.
%------------------------------------------------------------------*
\item Use that CW decomposition to compute the homology groups in all
dimensions.
%------------------------------------------------------------------*
\item Now consider $X$ to be the connected sum of the torus and the
projective plane.  
%%%%%%%%%%%%%%%%%%%%%%%%%%%%%%%%%%%%%%%%%%%%%%%%%%%%%%%%%%%%%%%%%%%
\begin{rlist}
%------------------------------------------------------------------*
\item Using this decomposition, compute all the homology groups of $X$
using the Mayer-Vietoris theorem.
%------------------------------------------------------------------*
\item Using this decomposition, give a presentation of the fundamental
group of $X$ using Van Kampen's theorem.
%------------------------------------------------------------------*
\end{rlist}
%%%%%%%%%%%%%%%%%%%%%%%%%%%%%%%%%%%%%%%%%%%%%%%%%%%%%%%%%%%%%%%%%%%
%------------------------------------------------------------------*
\end{alist}
%%%%%%%%%%%%%%%%%%%%%%%%%%%%%%%%%%%%%%%%%%%%%%%%%%%%%%%%%%%%%%%%%%%

\soln{jac -- Mon Aug 4 15:36:43 2003 -- A CW decomposition: look at
two projective planes $P_1$, $P_2$, each with one hole in them, and
one projective plane $P_3$ with two holes in it. The single-holed
planes are decomposed into two faces with four sides each. The
two-holed plane has two faces with six sides each.  These sides have
some identifications, described as follows: if $P_{11}$, $P_{12}$,
$P_{21}$, $P_{22}$, $P_{31}$, and $P_{32}$ are the distinct faces,
then $P_{11a}$, $P_{11b}$, $P_{11c}$, $P_{11d}$ are the edges of the
first face, then $P_{12a}$, $P_{12b}$, $P_{11c}$, $P_{11e}$ are the
edges of the second face.}
